\scanpage{21}
correspondante, \deleted{\unsure{qui|que}} de sorte qu'on aura un iso can.
\margin{plutôt, une \uline{équivalence}}
\[
(23)\qquad\mathcal M_{\unsure{G}}\simeq\prod_\alpha\mathcal M_{g_\alpha,\widehat I_\alpha}.
\]
Notons que $\Gamma=\operatorname{Aut}G$ opère sur $\mathcal M_G$ de façon évidente, et aussi sur le système $(\widehat I_\alpha,g_\alpha)_{\alpha\in S}$ donc sur le deuxième mb de (23), qui est un iso. compatible aux op. de $\Gamma$.

On peut considérer la multiplicité mixte
\[
(24)\qquad(\mathcal M_G,\Gamma)=\mathcal M_{[G]},
\]
qui a priori ne dépend que du type d'isomorphie $[G]$ de $G$, et qui s'interprète, en termes de ce dernier, comme la multiplicité schématique modulaire pour le groupoïde fibré (sur la $2$-cat. des multiplicités schématiques) des courbes standart relatives de type $[G]$.
\mspage{20}
